\textbf{结论名称：}角平分线定理（外角形式）

\textbf{适用条件：}在任意三角形中，一个外角的平分线与对边所在直线相交，外分对边。

\textbf{变量说明：}在 $\triangle ABC$ 中，$AD$ 是 $\angle A$ 的外角平分线，交边 $BC$ 的延长线于点 $D$。设 $AB=c$，$AC=b$，$BC=a$，则 $D$ 外分 $BC$，$BD$、$DC$ 表示线段长度（有向形式可约定方向）。

\textbf{核心公式：}
\[
\boxed{\frac{BD}{DC} = \frac{AB}{AC}}
\]

\textbf{使用场景：}求解与外角平分线相关的线段比或线段长；在竞赛题和模拟题中可快速转化边角关系，减少运算。

\textbf{简单例子：}在 $\triangle ABC$ 中，$AB=6$，$AC=4$，$BC=5$。作 $\angle A$ 的外角平分线交 $BC$ 延长线于点 $D$。由定理得 $\frac{BD}{DC}=\frac{6}{4}=\frac{3}{2}$。因 $AB>AC$，知 $D$ 位于 $BC$ 延长线上且顺序为 $B$—$C$—$D$。设 $DC=x$，则 $BD=5+x$，代入 $\frac{5+x}{x}=\frac{3}{2}$，解得 $x=10$。故 $DC=10$，$BD=15$。